% ============================================================
% code_and_data.tex
% Appendix A: Code and Data
% ============================================================

\label{sec:code_and_data}

This appendix documents what is required to reproduce the numerical
results reported in the paper.
The framework is implemented entirely in Python with a small set of
mainstream scientific dependencies; no specialised hardware is required
for any single experiment, although several long-running scans benefit
from leaving a workstation undisturbed for hours to days.

\subsection*{Repository and archived release}

\begin{itemize}
\item Source code, paper sources, and data files:\\
\url{https://github.com/JackDMenendez/discrete-causal-lattice}
\item Archived release (this version):\\
\url{https://doi.org/10.5281/zenodo.NNNNNNNN}
\item Citation metadata: \texttt{CITATION.cff} at the repository root.
\item License: MIT.
\end{itemize}

\subsection*{Software environment}

Tested on Windows 11 (MSYS2 UCRT64 shell) and Linux. The build
expects:

\begin{itemize}
\item Python~3 with the packages pinned in
\nolinkurl{virtual-env-requirements.txt}
(\texttt{numpy}~2.4.3, \texttt{scipy}~1.17.1,
\texttt{matplotlib}~3.10.8, \texttt{pandas}~3.0.1,
\texttt{networkx}~3.6.1, \texttt{seaborn}~0.13.2; full list in
the file).
\item GNU Make $\geq$ 4.3 (the grouped-target \texttt{\&:} syntax is used).
\item \texttt{pdflatex} and \texttt{bibtex} (MiKTeX or TeX Live) for
building this paper.
\end{itemize}

A clean build proceeds via the top-level \texttt{makefile}:

\begin{quote}
\texttt{make env}\hfill\textit{create the Python virtualenv}\\
\texttt{make tests}\hfill\textit{run the \texttt{pytest} suite under
\texttt{tests/}}\\
\texttt{make paper}\hfill\textit{compile \texttt{paper/main.tex} (3-pass
\texttt{pdflatex} + \texttt{bibtex})}
\end{quote}

\subsection*{Determinism and random-number use}

The framework is deterministic by construction: the tick rule is a
fixed map on the spinor amplitude, and orbital evolution depends
only on the initial state and the lattice geometry.
Random numbers enter in two places only:

\begin{enumerate}
\item \nolinkurl{src/core/TickScheduler.py} uses
\nolinkurl{np.random.default_rng(seed=42)} when an experiment requests
randomised tick ordering. \texttt{exp\_05} relies on this for the
shuffle-invariance test; all other experiments use the deterministic
sequential schedule.

\item \nolinkurl{src/experiments/exp_04_decoherence.py} draws
random phase scrambles via \nolinkurl{np.random.default_rng(42)}
for the many-bath ensemble; the per-realisation seed is
\texttt{i + 100} for the $i$-th member of the ensemble.
\end{enumerate}

With these seeds fixed at the values committed to the repository,
every experiment in this paper reproduces bit-for-bit on the pinned
software stack.

\subsection*{Running the numerical experiments}

Individual experiments are exposed as targets in
\nolinkurl{src/experiments/makefile}:

\begin{quote}
\texttt{make -C src/experiments exp\_12}\\
\texttt{make -C src/experiments exp\_16}
\end{quote}

Each target depends on the Python script of the same name, runs it
under the project virtualenv, and writes its log to
\texttt{data/exp\_NN\_*.txt} and its raw arrays to
\texttt{data/exp\_NN\_*.npy}. The longer scans
(\texttt{exp\_11}, \texttt{exp\_18}, \texttt{exp\_20}) have wall-clock
times of hours to days on consumer hardware; the per-experiment
documentation in \texttt{src/experiments/*.md} reports the runtime
observed during preparation of this paper.

\subsection*{Key parameter values}

The numerical experiments share a small set of physical and
discretisation parameters. Values used to produce the results
reported in the audit table (Table~\ref{tab:audit}) are summarised
in Table~\ref{tab:experiment_parameters}.

\begin{table}[htbp]
\centering
\footnotesize
\setlength{\tabcolsep}{4pt}
\caption{\small Parameter values for the numerical experiments cited
in the audit table. Lattice units throughout: distance in nodes,
time in ticks, $\hbar = c = 1$. \texttt{GRID} is the side length of
a cubic grid; \texttt{TICKS} is the total simulated time per run.}
\label{tab:experiment_parameters}
\begin{tabular}{@{}p{2.0cm}p{1.0cm}p{2.1cm}p{4.4cm}p{4.3cm}@{}}
\hline
\textbf{Experiment} & \textbf{GRID} & \textbf{TICKS} &
\textbf{Key parameters} & \textbf{Source} \\
\hline
\texttt{exp\_00} & $51^3$ & $20$ &
$\omega \in \{0,\, 0.5,\, 1.0\}$ &
\nolinkurl{exp_00_causal_cone.py} \\
\texttt{exp\_02} & $65^3$ & $4000$ &
test particle, mass well at centre &
\nolinkurl{exp_02_gravity_clock_density.py} \\
\texttt{exp\_06} & --- & $N \in \{4,\dots,64\}$ &
exact path-count integers &
\nolinkurl{exp_06_path_counting.py} \\
\texttt{exp\_08} & $65^3$ & $1500$ &
EM vs.\ gravity, Peierls $A\!\cdot\!v$ coupling &
\nolinkurl{exp_08_vacuum_twist.py} \\
\texttt{exp\_09} & $65^3$ & $4000$ &
$f\!\in\!(0,1)$, $\omega\!\in\!(0,2\pi)$ harmonic scan &
\nolinkurl{exp_09c_harmonic_hires.py} \\
\texttt{exp\_10} & $65^3$ & $8000$ &
$\omega = 0.1019$, \texttt{STRENGTH} $= 30.0$ &
\nolinkurl{exp_10_hydrogen_spectrum.py} \\
\texttt{exp\_11} & $65^3$ & $8000$ &
$\omega = 0.1019$, \texttt{STRENGTH} $= 30.0$,
\texttt{BURN\_IN} $= 500$ &
\nolinkurl{exp_11_quantization_scan.py} \\
\texttt{exp\_12} & $65^3$ &
$1500$ orbit / $2500$ scan &
$\omega_e = 0.1019$, $\omega_p = \pi/2$,
\texttt{STRENGTH} $= 30.0$ &
\nolinkurl{exp_12_hydrogen_twobody.py} \\
\texttt{exp\_16} & $65^3$ & $20000$ &
$\omega_p \in [0.3,\, \pi/2]$, \texttt{BURN\_IN} $= 4000$ &
\nolinkurl{exp_16_proton_mass_sweep.py} \\
\hline
\end{tabular}
\end{table}

The shared physical constants are: instruction frequency
$\omega_e = 0.1019$ (electron); proton instruction frequency
$\omega_p = \pi/2$ (maximum lattice mass); Coulomb strength
$30.0$; softening $0.5$; reference Bohr radius
$R_1 = 10.3$ nodes (the value at which $\omega \cdot R_1 = \pi/3$
saturates within $0.23\%$).

\subsection*{Data products}

Each experiment writes its raw arrays and a human-readable log to
\texttt{data/} at the repository root. The most consequential output
files are:

\begin{itemize}
\item \nolinkurl{data/exp_06_path_counts.npy} --- exact path-count
integers $P(N,a,b,c)$ underlying the discrete-correction prediction
P1 (Section~\ref{sec:predictions}).
\item \nolinkurl{data/exp_09_harmonic_hires.npy} ---
$(f,\omega)$-resolved spectral power, the source array for the
harmonic landscape (Figure~\ref{fig:harmonic_pair}).
\item \nolinkurl{data/exp_12_twobody_scan.npy} --- two-body
$k$-scan with the live proton, source for the
$k_\text{min} = 0.0970$ vs.\ $k_\text{Bohr} = 0.0971$ result
(Section~\ref{sec:hydrogen}, Figure~\ref{fig:exp12_twobody}).
\item \nolinkurl{data/exp_16_proton_mass_sweep.npy} --- proton-mass
sweep establishing the three-regime structure
(Section~\ref{sec:hydrogen}).
\end{itemize}

A representative output snippet from \texttt{exp\_12} follows;
the full log is at
\nolinkurl{data/exp_12_hydrogen_twobody_log.txt}.

\begin{quote}\small\ttfamily
[Test 3] Mini k-scan (8 values, 2500 ticks each)\\
\hphantom{xx}k\_Bohr\_fixed = 0.09709\\
\hphantom{xx}\dots\\
Two-body k\_min = 0.0970\quad vs fixed-well 0.0900\quad vs
k\_Bohr\_fixed 0.0971
\end{quote}

This is the four-significant-figure recovery of the Bohr radius
cited throughout the paper, produced with no parameter tuned to the
result.

\subsection*{Reproducing this paper}

The LaTeX source compiles from \texttt{paper/main.tex} via
\texttt{paper/build.cmd} (Windows) or \texttt{paper/build.sh} (POSIX),
or from the repository root via \texttt{make paper}. The build
performs three \texttt{pdflatex} passes around one \texttt{bibtex}
pass to settle cross-references and the bibliography.
The figures referenced in this paper are committed to
\texttt{figures/} as PDFs and \texttt{.tex} caption fragments;
re-running the experiments is not required to rebuild the document.
